\section{Faithfulness of the cancellation parameter}
\label{sec:cancellation-parameter-faithfulness}

For fixed $(m,r)$, the inverse polynomial \eqref{eq:cancellation-resolvent}
below is independent of the cancellation root $q$.  The root instead selects
the unique nonzero projective branch over $P=0$ which belongs to the affine
reconstruction open.  This section proves that the finite normalization
together with that open remembers the selection, even after stabilization.
The proof uses the reconstruction pole only through the cancellation
reconstruction residue introduced below.

For a normalized cancellation jet write
\[
 h_q(A)=q+h_1A+\cdots+h_rA^r
\]
and denote the resulting map by $F_q$.  Recall
\[
 A=1+xy^m,\qquad
 B=A^{r+1}z+y^{m+1}h_q(A),\qquad P=AB,qquad Q=y+xB.
                                                               \tag{6.1}
\]
The resolvent is
\begin{equation}
 \Psi_{P,Q,R}(T)
 =C\int_0^T\{1-t(Q-Pt)^m\}^r\,dt-R.                 \tag{6.2}
 \label{eq:cancellation-resolvent}
\end{equation}

\paragraph{Cancellation reconstruction residue.}
The affine component $A=0$ is intrinsically distinguished from $B=0$ by its
residue degree over $P=0$.  The function $y$ is a unit at its generic point,
so the marked reconstruction open carries the canonical quantity
\[
 \operatorname{CRR}(F_q)
 :=\left.\frac{B}{y^{m+1}}\right|_{A=0}=q.              \tag{6.2a}
 \label{eq:cancellation-reconstruction-residue}
\]
We call this the \emph{cancellation reconstruction residue}.  On the inverse
chart $D=A^{-1}$, where $B=PD$ and $y=Q-sP$, it is equivalently
\[
 \operatorname{CRR}(F_q)
 =\operatorname{res}_{A=0}\frac{PD}{(Q-sP)^{m+1}}.
\]
The proof below first recovers $A$ and $B$, normalizes $A$ using its unique
nondomain fiber, recovers the common weighted scaling of $y$ and $B$, and
then invokes the weight-zero residue
\eqref{eq:cancellation-reconstruction-residue}.

We first record the part of the target-boundary calculation which remains
valid in the presence of arbitrary stable variables.

\begin{lemma}[stable boundary-plane scaling]
\label{lem:stable-boundary-plane-scaling}
Let $\boldsymbol\tau$ be any finite tuple of variables.  Suppose a polynomial
automorphism of
\[
 \A^3_{P,Q,R}\times\A^{|\boldsymbol\tau|}
\]
preserves the two labelled cancellation boundary divisors $P=0$ and
$\Delta_{m,r}=0$.  Then for some $a,u\in k^*$ its ring pullback satisfies
\[
 \beta^*P=aP,
 \qquad
 \beta^*Q\equiv uQ\pmod P,
 \qquad
 \beta^*R\equiv u^{-m}R\pmod P.                       \tag{6.3}
 \label{eq:stable-boundary-plane-scaling}
\]
The conclusion permits arbitrary mixing of the stable variables.
\end{lemma}

\begin{proof}
Preservation of the prime divisor $P=0$ gives $\beta^*P=aP$, since the
only units of the polynomial ring are constants.  On the boundary plane,
Proposition~\ref{prop:thick-trace} gives the unconditional scheme-theoretic
intersection
\[
 Q^M G=0,
 \qquad
 M=mr(m+1),
 \qquad
 G=(r+1)RQ^m-C.                                        \tag{6.4}
\]
After adjoining $\boldsymbol\tau$, the two reduced components have coordinate
rings
\[
 k[R,\boldsymbol\tau]
 \quad\hbox{and}\quad
 k[Q,Q^{-1},\boldsymbol\tau].                          \tag{6.5}
\]
Their unit groups modulo $k^*$ have ranks zero and one, respectively, so an
automorphism cannot exchange them.  The induced automorphism $\bar\beta$ of
$k[Q,R,\boldsymbol\tau]$ therefore preserves the prime ideals $(Q)$ and
$(G)$.  Hence
\[
 \bar\beta(Q)=uQ,
 \qquad
 \bar\beta(G)=vG                                      \tag{6.6}
\]
for $u,v\in k^*$.  Substitution of the first equality into the second gives
\[
 (r+1)u^m\bar\beta(R)Q^m-C
 =v\{(r+1)RQ^m-C\}.
\]
Reduction modulo $Q$ gives $v=1$, and cancellation of $Q^m$ then gives
$\bar\beta(R)=u^{-m}R$.  This proves
\eqref{eq:stable-boundary-plane-scaling}.
\end{proof}

\begin{theorem}[stable cancellation-parameter faithfulness]
\label{thm:cancellation-parameter-faithfulness}
Let $k$ be an algebraically closed field of characteristic zero, fix
$m,r\geq1$, and let $q,q'$ be roots of $M_{m,r}$.  If the normalized
cancellation maps $F_q$ and $F_{q'}$ are stably polynomially left--right
equivalent, then
\[
 q=q'\qquad\hbox{and}\qquad h_q=h_{q'}.
                                                               \tag{6.7}
\]
Equivalently, the common finite normalization of
\eqref{eq:cancellation-resolvent}, together with
its distinguished affine reconstruction open and its map to the target,
faithfully remembers the selected cancellation root.
\end{theorem}

\begin{proof}
The argument has four intrinsic steps.

\smallskip
\noindent\emph{Recover the two affine factors.}
The case $mr=1$ has only one parameter root, so assume $mr>1$.  Then $P=0$
is the intrinsically labelled unramified target-boundary vertex.  After
adjoining the same tuple of identity variables to both maps, write the ring
identity supplied by a left--right equivalence as
\[
 \alpha^*F_{q'}^*=F_q^*\beta^*.                       \tag{6.8}
 \label{eq:parameter-faithfulness-equivalence}
\]
Lemma~\ref{lem:stable-boundary-plane-scaling} gives
\[
 \beta^*P=aP,
 \qquad
 \beta^*Q=uQ+P\varphi                                 \tag{6.9}
 \label{eq:parameter-faithfulness-target}
\]
for $a,u\in k^*$ and a polynomial $\varphi$, which may involve the stable
variables.

In the affine reconstruction open, the inverse image of $P=0$ has exactly
the two prime components $A=0$ and $B=0$.  Their generic residue degrees over
$P=0$ are respectively $1$ and $r+1$, by the projective boundary chart
\eqref{eq:cancellation-projective-branches}.  These labels survive
stabilization, so \eqref{eq:parameter-faithfulness-equivalence} cannot
exchange the two components.  Unique factorization and
\eqref{eq:parameter-faithfulness-target} therefore give a constant
$c\in k^*$ such that
\[
 \alpha^*A'=cA,
 \qquad
 \alpha^*B'=\frac acB.                                \tag{6.10}
 \label{eq:parameter-faithfulness-factors}
\]

\smallskip
\noindent\emph{Normalize $A$.}
The scalar $c$ is forced to be one.  Indeed, in the stabilized source the
fiber $A=1$ has coordinate ring
\[
 k[x,y,z,\boldsymbol\xi]/(xy^m),                       \tag{6.11}
\]
which is not a domain, whereas for every $t\ne1$ the fiber $A=t$ has
coordinate ring
\[
 k[y,y^{-1},z,\boldsymbol\xi]                          \tag{6.12}
\]
after solving $x=(t-1)y^{-m}$, and is a domain.  The pullback of the unique
nondomain fiber $A'=1$ under
\eqref{eq:parameter-faithfulness-factors} is $A=c^{-1}$; hence $c=1$.

\smallskip
\noindent\emph{Recover the weighted scaling.}
Put
\[
 X=\alpha^*x',\qquad Y=\alpha^*y'.
\]
Since $\alpha^*A'=A$, equation $A'-1=x'(y')^m$ becomes
\[
 XY^m=xy^m.                                             \tag{6.13}
 \label{eq:parameter-faithfulness-ufd}
\]
The stabilized source ring is a UFD, and $X,Y$ are nonassociate primes.
If $m>1$, comparison of prime multiplicities in
\eqref{eq:parameter-faithfulness-ufd} gives
\[
 X=\lambda^{-m}x,qquad Y=\lambda y                   \tag{6.14}
\]
for some $\lambda\in k^*$.  If $m=1$, unique factorization also permits the
formal interchange $X\sim y$, $Y\sim x$.  It cannot occur: reducing the
$Q$-identity in \eqref{eq:parameter-faithfulness-equivalence} modulo $B$
gives $Y\equiv uy\pmod B$, while
\[
 (B)\cap k[x,y,\boldsymbol\xi]=0                       \tag{6.15}
\]
because $B=A^{r+1}z+y^{m+1}h_q(A)$ has positive $z$-degree.  Thus the
interchange would make a nonzero linear combination of $x$ and $y$ belong to
$(B)$.  Consequently the scaling in (6.14) holds for every $m$.  The same
reduction modulo $B$ now gives $\lambda=u$.

Using \eqref{eq:parameter-faithfulness-factors} with $c=1$, the full
$Q$-identity is
\[
 uy+a u^{-m}xB
 =u(y+xB)+AB\,F_q^*(\varphi).                          \tag{6.16}
 \label{eq:parameter-faithfulness-Q}
\]
Cancel $B$.  If $a u^{-m}-u$ were nonzero,
\eqref{eq:parameter-faithfulness-Q} would imply
$A\mid x$, which is false.  Therefore
\[
 a=u^{m+1}.                                             \tag{6.17}
\]

\smallskip
\noindent\emph{Read the reconstruction residue.}
At the affine prime $A=0$, apply
\eqref{eq:cancellation-reconstruction-residue}.
Equations \eqref{eq:parameter-faithfulness-factors} and (6.14)--(6.17)
carry the analogous residue for $F_{q'}$ to
\[
 \left.\frac{u^{m+1}B}{(uy)^{m+1}}\right|_{A=0}=q.
\]
Thus $\operatorname{CRR}(F_{q'})=\operatorname{CRR}(F_q)$ and $q'=q$.  The
finite cancellation recursion gives a unique normalized
degree-$r$ jet over each parameter root, so $h_{q'}=h_q$ as well.
\end{proof}

\begin{remark}[cancellation reconstruction residue]
The cancellation reconstruction residue is the first coefficient after the
canonical leading pole
in
\[
 PD^{r+2}-z
 =(Q-sP)^{m+1}D^{r+1}
 \{q+h_1D^{-1}+\cdots+h_rD^{-r}\}.                     \tag{6.19}
\]
The four-step proof is precisely the intrinsic normalization needed before
this Laurent coefficient can be compared.
\end{remark}

\begin{corollary}[distinct roots give distinct stable classes]
\label{cor:cancellation-root-classes}
For fixed $(m,r)$, the $mr$ geometric roots of the squarefree polynomial
$M_{m,r}$ give $mr$ pairwise stably polynomially left--right inequivalent
cancellation maps.  Cancellation maps belonging to different pairs $(m,r)$
of the same inverse degree remain separated by their boundary signatures.
Consequently, if $n=N-1$, the cancellation branches together with one
weighted class give the degreewise lower bound
\[
 1+\sum_{\substack{r\mid n\\r<n}}(n-r)
 =1+n\tau(n)-\sigma(n).
\]
\end{corollary}
